% ===========================================================================
% RELATED WORKS
% File: sections/02_related_works.tex
% Built sub-section by sub-section
% ===========================================================================

\section{Related Works}
\label{sec:related}

% ---------------------------------------------------------------------------
\subsection{Agrivoltaic Systems: Architecture, Crop Response, and
Land-Use Efficiency}
% ---------------------------------------------------------------------------

\hspace{1cm}The concept of co-locating solar panels and crops on the same
land parcel was first proposed by \citet{goetzberger1982}, who noted that
partial shade from PV modules could reduce transpiration stress while the
modules themselves benefited from evaporative cooling. The idea remained
largely theoretical until \citet{dupraz2011} published the first rigorous
field experiment and introduced the Land Equivalent Ratio (LER) as the
standard metric for the dual land-use advantage. Measuring durum wheat
yield and PV output simultaneously in southern France, they demonstrated
LER values of 1.35--1.73 depending on shade level, establishing the
foundational result that the combined system outperforms two separate
monocultures by 35--73\% on the same land area.

Subsequent field studies confirmed that the benefit extends beyond
electricity and yield. \citet{marrou2013} demonstrated that the microclimate
under APV panels --- characterised by soil temperature reductions of
1.3--2.3\,$^{\circ}$C and reduced peak radiation load --- can partially
offset or reverse the yield penalty of reduced PAR for shade-adapted
lettuce, cucumber, and wheat. In semi-arid Arizona, \citet{barron2019}
showed that agrivoltaics simultaneously increased jalapeno yield by
3$\times$, reduced water consumption by 65\%, and improved panel
efficiency by 1--9 percentage points through evaporative cooling,
providing the first comprehensive food-energy-water nexus evidence for
the technology.

The integration of bifacial modules --- which harvest both direct
front-side and diffuse rear-side irradiance --- represents a significant
performance advance. \citet{riaz2022} developed a Light-Productivity-Factor
(LPF) framework that accounts for the PAR saturation threshold of the
specific crop. Applied to bifacial tomato production in Arizona, this
approach yielded the highest published two-component LER of 1.94,
demonstrating that crop-specific geometry optimization can push the
land-use advantage substantially beyond geometry-agnostic designs. The
PAR saturation concept introduced by \citet{riaz2022} --- the observation
that above a crop-specific irradiance threshold, additional light generates
no additional photosynthesis and the surplus can be productively captured
as electricity --- forms the theoretical basis for the crop yield model
adopted in the present study.

\citet{mazzeo2025} extended this approach to open-field lettuce in Italy,
achieving LER = 1.85 under optimised panel density and identifying the
circular reference problem in LER$_{\text{PV}}$ calculation: using the
same APV array as both system and reference inflates the LER value
artificially. This issue motivates the independent dense-pack PV reference
(Definition~A) adopted here. At the system scale, a meta-analysis of 127
agrivoltaic case studies by \citet{semeraro2024} found a mean LER of
$1.34 \pm 0.18$, with higher values consistently associated with bifacial
modules, shade-tolerant crops, and optimised row spacing. Water savings
from APV shading, quantified by \citet{elamri2018} across a lettuce trial
in southern France, reached 14--29\% of growing-season ET$_0$ --- a range
consistent with the corrected values reported in the present study.

Despite this body of evidence, all cited agrivoltaic studies share a common
limitation: the crop residue stream is treated as waste or a soil amendment,
severing the link between the agrivoltaic system and the biogas economy that
the same semi-arid farming regions are increasingly participating in.

\subsection{Anaerobic Digestion: Fundamentals, Feedstocks, and
Thermal Management}
% ---------------------------------------------------------------------------
 
\hspace{1cm}Anaerobic digestion is a multi-stage biological process in
which consortia of microorganisms decompose organic matter in the absence
of oxygen, producing biogas --- typically 55--70\% methane and 30--45\%
CO$_2$ --- and a nutrient-rich liquid digestate~\citep{angelidaki_defining_2009,mata-alvarez_critical_2014}. The four sequential stages --- hydrolysis, acidogenesis,
acetogenesis, and methanogenesis --- each involve distinct microbial
communities with different kinetic sensitivities to temperature, pH, and
substrate composition~\citep{alengebawy2024}. This sequential dependency
makes methanogenesis the rate-limiting step and temperature the dominant
operational parameter: a 10\,$^{\circ}$C decrease in digester temperature
below the mesophilic optimum (37\,$^{\circ}$C) can suppress methane yield
by 30--50\%, following an Arrhenius-type relationship~\citep{lindorfer2008}.
 
The biochemical methane potential (BMP) of a substrate quantifies the
maximum volume of methane recoverable per unit of volatile solids (VS)
under controlled conditions. \citet{lallement2023} reported BMP values
of 200--280\,NmL\,CH$_4$\,gVS$^{-1}$ for cattle manure and
250--350\,NmL\,CH$_4$\,gVS$^{-1}$ for tomato plant residues --- ranges
that bracket the reference value of 300\,NmL\,CH$_4$\,gVS$^{-1}$ at
37\,$^{\circ}$C adopted in the present study. \citet{feng2013} provided
a five-point temperature response curve for tomato plant residues from
20\,$^{\circ}$C to 40\,$^{\circ}$C, demonstrating that BMP drops from
300\,NmL\,gVS$^{-1}$ at 37\,$^{\circ}$C to approximately
195\,NmL\,gVS$^{-1}$ at 20\,$^{\circ}$C --- a 35\% suppression that is
directly relevant to cold-season operation at Konya and Freiburg. This
dataset provides the empirical basis for the V2 model validation reported
in Section~\ref{sec:results}. \citet{pilarski2025} extended this analysis
to cereal straw and vegetable mixtures, confirming the Arrhenius coefficient
$\theta = 0.069$\,$^{\circ}$C$^{-1}$ used throughout the present study and
reporting that co-digestion of crop residues with cattle manure at a 1:3
VS ratio increases BMP by 12--18\% through synergistic microbial diversity.
 
The thermal energy demand of a continuously stirred tank reactor (CSTR)
digester operating under mesophilic conditions in a temperate climate is
substantial. \citet{lindorfer2008} showed that maintaining a 37\,$^{\circ}$C
digester against a 0\,$^{\circ}$C ambient requires heat input equivalent
to 15--30\% of gross biogas energy, depending on insulation quality and
substrate inlet temperature. This self-consumption directly reduces the
net energy yield and economic return of the plant. Solar thermal
collectors~\citep{temiz2024} and PV-powered electrical resistance
heating~\citep{pal2024} have both been proposed as alternatives to
biogas self-consumption. \citet{pal2024} demonstrated that replacing
biogas self-heating with PV electricity in a grid-connected configuration
can improve the NPV of a biogas plant by up to 8.7$\times$ compared to
a non-heated reference, by redirecting the saved biogas to higher-value
grid injection. Their analysis, however, treats the PV array as a
separate installation without modelling the agrivoltaic geometry or crop
yield effects.
 
On the feedstock supply side, the strategic importance of agricultural
biogas and biomethane is reflected in accelerating policy commitments.
The EU REPowerEU Plan targets 35\,billion m$^3$ of sustainable biomethane
per year by 2030~\citep{ec2022}, requiring a nine-fold increase from 2022
levels. Germany's EEG 2023 and the EU RED III mandate feed-in premiums
for biomethane grid injection, creating a market price approximately
5.5 times the wholesale electricity price in Central European
markets~\citep{nikzad2026, eba2025, rabobank2025}. In Sub-Saharan Africa,
where Burkina Faso's national grid relies on 63\% diesel generation, the
IEA projects biogas from agricultural residues as a priority pathway for
rural electrification and LPG substitution by 2030~\citep{iea2025}.
These policy contexts directly motivate the site-specific biomethane
premium scenarios (S5) and carbon credit scenarios (S6) evaluated in the
present study.

% ---------------------------------------------------------------------------
\subsection{PV--Biogas Hybrid Systems: Architecture and Economics}
% ---------------------------------------------------------------------------

\hspace{1cm}The complementarity between solar PV and biogas is
thermodynamic as much as economic. PV generation is intermittent and
peaks in summer, while biogas is dispatchable and its heating demand
peaks in winter --- a seasonal anti-correlation that makes the two
technologies natural partners in a hybrid energy
system~\citep{pal2024, temiz2024}. Early hybrid configurations
focused on off-grid rural electrification, combining a biogas generator
for baseload with PV for daytime peak demand and battery storage for
overnight supply. More recently, grid-connected configurations that
exploit feed-in premiums for both electricity and biomethane have
become an important economic model in European
markets~\citep{eba2025, nikzad2026}.

Techno-economic analyses of PV--biogas hybrids consistently report
that the optimal configuration depends critically on the local
electricity tariff structure and the relative prices of grid
electricity, LPG, and biomethane. \citet{pal2024} showed that in
markets where grid electricity is cheap relative to biogas, using PV
electricity to heat the digester (rather than burning biogas for
self-heating) frees the entire biogas output for sale, improving IRR
by 3--8 percentage points. \citet{temiz2024} demonstrated that a
heat pump powered by APV electricity can maintain thermophilic
digestion (52\,$^{\circ}$C) year-round in a continental climate,
achieving full thermal self-sufficiency with a coefficient of
performance (COP) of 3.2--4.1, equivalent to a 70\% reduction in
heating energy relative to direct resistance heating.

Multi-generation system analyses --- simultaneously producing
electricity, heat, biomethane, and digestate fertilizer --- have
been explored by \citet{roldanporta2023}, who reported that joint
optimization of all four output streams increases system NPV by
22--34\% compared to single-output optimization. Their framework,
however, does not include a land-use efficiency metric or crop yield
modelling, treating the agricultural component only as a feedstock
source.

A critical gap in all cited PV--biogas studies is the absence of
agrivoltaic geometry modelling: the PV array is treated as an
independent installation with no interaction with the underlying
crop canopy. This simplification ignores the fact that in an
integrated APV--AD system, the panel geometry simultaneously
determines PV yield, crop PAR availability, residue production,
and VS input to the digester --- four interlinked variables that
must be co-optimized rather than treated sequentially.

% ---------------------------------------------------------------------------
\subsection{The Land Equivalent Ratio: History, Extensions,
and Limitations}
% ---------------------------------------------------------------------------

\hspace{1cm}The Land Equivalent Ratio was introduced by
\citet{mead1980} for intercropping systems, where it quantifies
the relative land area under separate monocultures that would be
required to produce the same combined yields as an intercropped
system. An LER of 1.4, for example, means that 40\% more land
would be needed under monoculture to achieve the same total
output --- a powerful argument for integration in land-constrained
contexts. \citet{dupraz2011} adapted the metric to agrivoltaics
by replacing one of the two intercrop outputs with PV electricity,
defining the two-component agrivoltaic LER as:
\[
    LER_{2C} = \frac{Y_{\text{crop,APV}}}{Y_{\text{crop,open}}}
             + \frac{E_{\text{PV,APV}}}{E_{\text{PV,ref}}}
\]
where the denominators are the yields of the same crop and PV
technology operated separately on equivalent land areas.

The choice of reference system for the PV denominator has emerged
as a methodological fault line in the agrivoltaic literature.
Three definitions are in common use, each yielding substantially
different numerical LER values for the same physical system.
The first --- used by \citet{dupraz2011} and most early studies
--- compares APV electricity output to a dense-pack ground-mount
array at the same tilt, giving a rigorous independent reference
(Definition~A in the present study). The second compares the APV
output to the same APV array itself, making LER$_{\text{PV}}$
equal to the Ground Coverage Ratio by definition --- a circular
reference identified and criticised by \citet{mazzeo2025}.
The third normalizes by actual irradiance received rather than
potential output, conflating system performance with resource
availability. \citet{mazzeo2025} showed that the choice between
Definition~A and the circular definition changes reported LER
values by 0.15--0.35, making cross-study comparisons unreliable
unless the reference is stated explicitly. The present study
reports results under both Definition~A (primary) and
Definition~B (circular, for literature comparability), and
flags the numerical difference explicitly in the Results section.

A second unresolved limitation of the classical LER framework is
its restriction to two outputs. In an integrated APV--AD system,
three distinct land-use outputs are produced simultaneously:
crop yield, PV electricity, and biogas energy. Collapsing the
third output into the economic analysis --- as a revenue item
in an NPV calculation --- obscures its contribution to land-use
efficiency and prevents direct comparison with classical LER
values. No prior study has proposed a three-component LER
formulation that preserves the dimensional consistency and
interpretive clarity of the original metric while accommodating
a biogas energy term. The eLER proposed in the present study
fills this gap, defining $LER_{\text{biogas}}$ as the ratio of
APV-system biogas output to the output of a standalone AD unit
processing the same land area under open-field conditions, and
summing all three components into a single dimensionless indicator
of multi-output land productivity.


 % ---------------------------------------------------------------------------
\subsection{The APV--AD Frontier: First Integration Studies}
% ---------------------------------------------------------------------------

\hspace{1cm}An early study to explicitly couple agrivoltaic and
anaerobic digestion optimization was \citet{nikzad2026}, who applied
a techno-economic framework to a grid-connected APV--AD plant in
northern Italy. Their system integrates single-axis tracking PV panels
above an agricultural field, with biogas produced from on-site biomass
feeding a combined heat and power (CHP) unit or a biomethane upgrading
system for grid injection. Using NPV as the objective function and land
occupation as a secondary constraint, they found that on-grid
configurations with heat pump integration require 17--20 times less
land than off-grid alternatives and achieve NPV values up to
\texteuro2.88 million at the optimised design. The biomethane grid
injection scenario, valued at 5.5 times the electricity purchase price
under EEG 2023 and RED III incentives, emerged as the economically
preferred valorisation strategy --- a finding that directly motivates
Scenario S5 of the present study.

Despite its pioneering contribution, the work of \citet{nikzad2026}
has four significant limitations that the present study addresses.
First, the analysis is geographically limited to a single
Mediterranean location (northern Italy, Cfb climate), preventing
any inference about how the APV--AD synergy varies with solar
resource, ambient temperature, or growing-season length.
Second, the crop beneath the panels is treated exclusively as a
biomass feedstock source: no crop yield model is implemented,
and the effect of panel shading on photosynthesis, PAR
distribution, or harvest quantity is not quantified.
Third, no land-use efficiency metric is defined beyond raw land
occupation; the eLER concept --- or any equivalent multi-output
indicator --- is absent, making it impossible to assess how
efficiently the integrated system uses land relative to three
separate monocultures.
Fourth, the optimization searches only the economic design space
(CHP vs. biomethane, on-grid vs. off-grid) without jointly
optimizing the APV geometry (tilt, row spacing, module height)
and AD design (digester volume, HRT, heating fraction)
simultaneously.

Table~\ref{tab:positioning} summarises the positioning of the
present study relative to the most relevant prior works across
the agrivoltaic, AD integration, and LER literatures.

\begin{table}[H]
\caption{Positioning of the present study relative to key prior works.
\checkmark~=~fully addressed; $\circ$~=~partially addressed;
$\times$~=~not addressed.}
\label{tab:positioning}
\centering
\small
\begin{tabular}{lcccccc}
\toprule
\textbf{Study} &
\textbf{Multi-} &
\textbf{Crop} &
\textbf{AD} &
\textbf{Joint} &
\textbf{LER} &
\textbf{4E} \\
& \textbf{climate} &
\textbf{yield} &
\textbf{thermal} &
\textbf{optim.} &
\textbf{metric} &
\textbf{analysis} \\
\midrule
\citet{dupraz2011}    & $\times$ & \checkmark & $\times$ & $\times$ & \checkmark & $\times$ \\
\citet{riaz2022}      & $\times$ & \checkmark & $\times$ & $\times$ & \checkmark & $\times$ \\
\citet{mazzeo2025}    & $\times$ & \checkmark & $\times$ & $\circ$  & \checkmark & $\times$ \\
\citet{zainali2023}   & $\circ$  & $\times$   & $\times$ & $\times$ & $\times$   & $\times$ \\
\citet{pal2024}       & $\times$ & $\times$   & \checkmark & $\times$ & $\times$ & $\circ$  \\
\citet{temiz2024}     & $\times$ & $\times$   & \checkmark & $\times$ & $\times$ & $\circ$  \\
\citet{nikzad2026}    & $\times$ & $\times$   & \checkmark & $\circ$  & $\times$ & $\circ$  \\
\textbf{This study}   & \checkmark & \checkmark & \checkmark & \checkmark & \checkmark & \checkmark \\
\bottomrule
\end{tabular}
\end{table}

% ---------------------------------------------------------------------------
\subsection{Multi-Climate Comparative Studies and Research Positioning}
% ---------------------------------------------------------------------------

\hspace{1cm}Multi-site comparative analyses in agrivoltaic research
remain rare. \citet{zainali2023} modelled shading factors and POA
irradiance across five European locations using a validated geometric
model, but without crop yield, economic, or water analysis. Their study
confirmed that the shading geometry is strongly latitude-dependent ---
the same panel configuration produces significantly different ground
irradiance distributions at 38°N (Spain) versus 52°N (Netherlands) ---
motivating the multi-latitude design of the present study. \citet{semeraro2024}
conducted a meta-analysis across 127 agrivoltaic case studies spanning
Europe, Asia, and North America, finding that climate zone explains
approximately 18\% of LER variance after controlling for crop species
and panel technology, but without isolating the mechanism.

No prior study has applied an integrated APV--AD framework to more than
one climate zone, nor has any study included a site from tropical
Sub-Saharan Africa (BSh tropical), where the combination of year-round
near-optimal digester temperatures (ambient T$_{\text{amb}} \geq$
22\,$^{\circ}$C throughout the year), high solar irradiance
(\SI{2050}{kWh\,m^{-2}}), and acute energy poverty creates a
particularly strong case for integrated APV--AD deployment. The inclusion of
Ouagadougou, Burkina Faso in the present study --- alongside Konya
(BSk), Almer\'{i}a (BSh hot), and Freiburg (Cfb) --- enables a quantitative comparison of APV--AD performance across the full
semi-arid-to-temperate spectrum.

The present study extends prior work by:
(i) applying a joint APV--AD optimizer across four climate zones,
(ii) introducing a three-component eLER as the optimization objective,
(iii) identifying and mechanistically explaining a non-monotone relationship
between annual solar resource and eLER through the PAR-saturation
mechanism, and
(iv) providing a 4E analysis --- Energetic, Economic, Energo-economic,
and Environmental --- for an integrated APV--AD system across
contrasting climate zones. Together, these contributions address
the six research gaps identified in Section~\ref{sec:intro} and
establish a replicable framework for the site-specific deployment
assessment of integrated APV--AD systems in food-energy-water nexus
planning.
